% Pregunta 274 (2018-enero-OEP2017, n.º 30): curva en J; tras la devaluación el saldo comercial cae por el efecto precio y luego sube por encima del inicial cuando se cumple Marshall-Lerner. Autor: Víctor Gutiérrez Marcos
\documentclass[tikz,border=6pt]{standalone}
\input{tikz-tcee}
\begin{document}
\begin{tikzpicture}[tcee]
  \ejes{$t$}{$SC$}{6.4}{5}
  % Saldo inicial y momento de la devaluación
  \draw[guia] (0,2.6) -- (6.2,2.6);
  \node[left] at (0,2.6) {$SC_0$};
  \draw[guia] (1.5,0) -- (1.5,4.8);
  \node[below] at (1.5,0) {$t_0$};
  \node[etiqueta,font=\footnotesize,align=center,anchor=south] at (1.5,4.8) {devaluación};
  % Sin Marshall-Lerner (opciones a y c): no se recupera
  \draw[tceegris,very thick,dashed] (1.5,2.6) .. controls (1.75,1.9) and (2.0,1.45) .. (2.5,1.25)
        .. controls (3.5,0.95) and (5.0,0.8) .. (6.0,0.7);
  \node[etiqueta,tceegris,font=\footnotesize,align=left,anchor=west] at (3.9,0.4) {sin Marshall-Lerner};
  % Curva en J
  \draw[curva1] (0.2,2.6) -- (1.5,2.6) .. controls (1.75,1.9) and (2.0,1.55) .. (2.4,1.55)
        .. controls (3.2,1.55) and (4.2,3.2) .. (6.0,4.1);
  \node[etiqueta,tceemorado,right] at (6.0,4.1) {$J$};
  \node[punto] at (2.4,1.55) {};
  % Rótulos de los dos efectos
  \node[etiqueta,font=\footnotesize,align=center,anchor=east] at (1.42,1.7) {efecto\\precio};
  \node[etiqueta,font=\footnotesize,align=center,anchor=west] at (2.0,4.1) {efecto cantidad\\(Marshall-Lerner)};
  \draw[flecha] (3.4,3.75) -- (4.1,3.1);
\end{tikzpicture}
\end{document}
